\scanpage{36}
% archive sheet number: 35
On trouve donc
\[
\sigma_0^{\unsure{\ell|\varepsilon}}=\mathrm{id}\quad\Longleftrightarrow\quad
1+\unsure{\eta_0|\gamma_0}+\cdots+\eta_0^{\ell-1}=0
\]
donc

Ordre de $\sigma_0=$ ordre de $\eta_0\in k^*$ \uline{sauf} si $\eta_0=1$ \inserted{i.e. $\det\sigma_{0V}=1$}, auquel cas $\sigma_{0V}$ est la transvection
$\begin{pmatrix}1&\alpha_0\\0&1\end{pmatrix}$
(qui est d'ordre $1$ ssi $\alpha_0=0$, sinon d'ordre $p$)
\[
\eta_0=1\quad\left\{\begin{aligned}
\text{a)}\quad&\alpha_0=1\quad\sigma_0=\mathrm{id}\qquad\text{\uline{cas impropre}}\\
\text{b)}\quad&\alpha_0\ne1\quad\left\{\begin{aligned}
p=0\quad&\sigma_{0V}\text{ d'ordre infini, a fortiori }\sigma_0\text{ d'o. inf.}\\
p>0\quad&\sigma_{0V}\text{ d'ordre }p,\quad\sigma_0^p=\deleted{\text{translation}}\inserted{\text{\unsure{au pt.}}}\\
&\deleted{\unsure{(n,0)}}\quad\left\{\begin{aligned}
&\deleted{p\mid n\ \Longrightarrow}\ \sigma_0^p=\mathrm{id}\\
&\deleted{p\nmid n\quad\sigma_0\text{ d'ordre }p^2}
\end{aligned}\right.
\end{aligned}\right.
\end{aligned}\right.
\]
\par\noindent\rule{.4\linewidth}{.4pt}\par
\[
\begin{aligned}
\sigma_1^2(x,y)&=(\lambda_1x+\lambda_1(1-\lambda_1)y,(1-\lambda_1)x+(\lambda_1+(1-\lambda_1)^2)y)\\
&=(\lambda_1x+\lambda_1(1-\lambda_1)y,(1-\lambda_1)x+(\lambda_1^2-\lambda_1+1)y)
\end{aligned}
\]
\[
\sigma_1=\begin{pmatrix}0&\lambda_1\\1&1-\lambda_1\end{pmatrix}
\]
\[
\det(t\,\mathrm{id}-\sigma_1)=t^2-(1-\lambda_1)t-\lambda_1
=\deleted{t^2-(1+\delta)t+\delta}
\]

Valeurs propres \deleted{\unsure{$\eta,\eta'$}}
\inserted{$1,-\lambda_1=\deleted{\delta}\eta_1$}
\[
\deleted{\left\{\begin{aligned}-\lambda_1&=\eta\eta'\\1-\lambda_1&=\eta+\eta'\end{aligned}\right.}
\qquad\deleted{-\lambda_1=\delta_1=\delta}
\]
\[
\deleted{\left\{\begin{aligned}\eta\eta'&=\delta\\\eta+\eta'&=1+\delta\end{aligned}\right.}
\]
% margin note written slanted
\margin{\unsure{alors|choix}}
\[
\deleted{\eta^{\unsure{n|m}}=\eta'{}^n=1}\quad\deleted{\Longrightarrow\delta^n=1}
\]
\[
\deleted{\frac{\eta}{\delta}}\qquad\deleted{\frac{\eta'}{\delta}}\qquad
\deleted{\frac{\eta}{\delta}}\qquad\deleted{\frac{\eta'}{\delta}}\qquad
\deleted{y=}\quad x=\lambda_1y
\]
\[
\deleted{\unsure{\delta\ne1}}
\]

Ordre de $\sigma_1=$ ordre de $\eta_1\in k^*$, \uline{sauf si $\delta=1$} i.e. $\lambda_1=-1$
\[
\sigma_1(x,y)=(-y,x+2y)\quad\text{transvection}
\]
\unsure{\deleted{non}|$\ne$} identité, d'ordre infini si $p=0$, d'ordre $p$ sinon.
